% =====================================================================
%  appendix.tex — נספחים: מילון מונחים + דף למקרר
% =====================================================================
\chapter*{נספח א׳: מילון מונחים קצר}
\addcontentsline{toc}{chapter}{נספח א׳: מילון מונחים קצר}
\markboth{מילון מונחים}{}
\thispagestyle{plain}

כמה מונחים שחוזרים לאורך המדריך, בהסבר קצר ובשפה פשוטה.

\newcommand{\term}[2]{\par\addvspace{5pt}{\sffamily\bfseries\color{navy}#1}\enspace{\color{coral}\textbf{·}}\enspace #2\par}

\term{חיזוק חיובי}{הוספה של משהו נעים, כמו חטיף, שבח או משחק, מיד אחרי התנהגות רצויה, כדי שהכלב
יבחר בה שוב. זו הגישה שמומלצת לאורך המדריך (פרק 1).}
\term{עומס חושי}{מצב שבו מצטברים יותר מדי גירויים, כמו רעש, אנשים, ריחות ותנועה, והכלב מתקשה להקשיב
או להירגע.}
\term{מרחק נוח}{המרחק מגירוי שבו הכלב מבחין בו אבל עדיין רגוע מספיק כדי להקשיב ולקחת חטיף. בעזרתו
עובדים על חשיפה הדרגתית (סולם החשיפה, פרק 1).}
\term{ברכיצפלי}{כלב בעל אף ופנים קצרים, כמו בולדוג או פאג. לכלבים כאלה סיכון גבוה יותר למכת חום
(פרק 2).}
\term{מכת חום}{התחממות־יתר מסוכנת של הגוף. היא מצב חירום שדורש פנייה מיידית לווטרינר (פרק 4).}
\term{קסיליטול}{ממתיק שנמצא בחלק מהמסטיקים ומהמוצרים ״ללא סוכר״ והוא מסוכן לכלבים (פרק 4).}
\term{שבב אלקטרוני}{שבב זיהוי קטן שמושתל מתחת לעור. וטרינר או השירות הווטרינרי העירוני יכולים לסרוק
אותו ולאתר את הבעלים.}
\term{שער כפול}{שני שערים ברצף בכניסה לגינת כלבים, שמונעים מכלבים לברוח החוצה (פרק 2).}
\term{ריחרוח}{הדרך העיקרית שבה כלבים אוספים מידע על העולם. טיול עם הרבה ריחרוח הוא גם העשרה
וגם דרך להירגע.}
\term{תוכנית ב׳}{אדם וזמן חלופיים שמוכנים מראש למקרה שהשגרה מתפרקת (פרק 3).}

\begin{fieldbox}
המילון קצר בכוונה. אם נתקלתם במונח שלא מופיע כאן, שאלו את הווטרינר או את המאלף: אין שאלה
טיפשית כשמדובר במי שגר אתכם בבית.
\end{fieldbox}
\clearpage

% ---------------------------------------------------------------------
\chapter*{נספח ב׳: דף למקרר}
\addcontentsline{toc}{chapter}{נספח ב׳: דף למקרר}
\markboth{דף למקרר}{}
\thispagestyle{plain}

כל המדריך בעמוד אחד. הדפיסו, תלו על המקרר, והשאירו ליד הרצועה.

\begin{tcbraster}[raster columns=1,raster row skip=14pt,raster before skip=10pt,
  raster every box/.style={fontupper=\large,top=7pt,bottom=7pt,before upper={\setlist[itemize]{itemsep=5pt,label={\color{coral}\small$\bullet$}}}}]
  \begin{tcolorbox}[enhanced,colback=sage!12!white,colframe=sagedark,arc=3mm,boxrule=0.9pt,
     title={\sffamily\bfseries\paw[white]\enspace לפני כל טיול},colbacktitle=sagedark,left=8pt,right=8pt]
    \begin{itemize}
      \item רצועה, קולר או רתמה תקינים, ותג עם טלפון עדכני.
      \item שקיות (ועוד אחת), מים וחטיפים קטנים.
      \item בדיקת מזג אוויר: בימים חמים, הנחת גב כף היד על המדרכה ובחירת מסלול מוצל.
      \item בחירה בין שלושת המסלולים הקבועים: קצר, מוצל או ארוך.
    \end{itemize}
  \end{tcolorbox}
  \begin{tcolorbox}[enhanced,colback=navy!6!white,colframe=navy,arc=3mm,boxrule=0.9pt,
     title={\sffamily\bfseries\infoicon\enspace בדרך},colbacktitle=navy,left=8pt,right=8pt]
    \begin{itemize}
      \item קוראים את הגוף: גוף רפוי זה סימן טוב, וזנב מקופל, קפיאה או ליקוק שפתיים הם סימנים לתת מרחק.
      \item מעלית: שואלים לפני שנכנסים, רצועה קצרה, ומחכים למעלית הבאה אם צריך.
      \item מפגש עם כלב אחר: שואלים את הבעלים, מתקרבים בקשת, ומסיימים אחרי כמה שניות.
      \item אוספים כל צואה ומקצרים את הרצועה ליד ילדים, רוכבים ובתי קפה.
    \end{itemize}
  \end{tcolorbox}
  \begin{tcolorbox}[enhanced,colback=alert!5!white,colframe=alert,arc=3mm,boxrule=1.2pt,
     title={\sffamily\bfseries\medicon\enspace כשמשהו משתבש},colbacktitle=alert,left=8pt,right=8pt]
    \begin{itemize}
      \item הכלב בלע משהו חשוד? מרחיקים אותו, מזהים מה זה, ומתקשרים לווטרינר. \textbf{לא גורמים להקאה בעצמכם.}
      \item התנשפות חזקה, חולשה, הקאות, קריסה או פרכוסים: נוסעים מיד לווטרינר.
      \item חום: עוברים לצל, מציעים מים וחוזרים הביתה.
      \item מספרי החירום נמצאים בכרטיס החירום (פרק 4).
    \end{itemize}
  \end{tcolorbox}
\end{tcbraster}

\vfill
\begin{center}
  {\sffamily\small\color{sagedark}\paw[sage]\enspace המדריך הלא-רשמי לחיים עם כלב בעיר הגדולה\enspace\paw[sage]}
\end{center}
\clearpage
